\chapter{Guía de arranque del sistema}
\label{cap:AnexoK}

Este anexo describe los pasos necesarios para poner en marcha el sistema en un entorno de desarrollo local. Toda la infraestructura se gestiona mediante Docker Compose y un \texttt{Makefile} que centraliza los comandos habituales.

\section{Requisitos previos}

\begin{itemize}
  \item \textbf{Docker} 24 o superior con el plugin \texttt{docker compose} (v2).
  \item \textbf{GNU Make}.
  \item \textbf{Git}.
  \item \textit{Opcional — modo GPU:} controlador NVIDIA y \texttt{nvidia-container-toolkit} instalados en el sistema anfitrión.
\end{itemize}

\section{Configuración del entorno}

Antes de arrancar por primera vez es necesario crear el fichero de variables de entorno a partir de la plantilla incluida en el repositorio:

\begin{lstlisting}[language=bash]
cp .env.example .env
\end{lstlisting}

Los valores por defecto son válidos para un entorno local de desarrollo. Las variables más relevantes se recogen en la tabla~\ref{tab:env-vars}.

\begin{table}[h]
\centering
\caption{Variables de entorno principales.}
\label{tab:env-vars}
\begin{tabular}{lll}
\hline
\textbf{Variable} & \textbf{Valor por defecto} & \textbf{Descripción} \\
\hline
\texttt{POSTGRES\_PASSWORD} & \texttt{changeme}   & Contraseña de PostgreSQL \\
\texttt{SECRET\_KEY\_BASE}  & (valor de ejemplo)  & Clave secreta de Phoenix \\
\texttt{HF\_TOKEN}          & \texttt{hf\_xxx...} & Token de Hugging Face (modo GPU) \\
\texttt{BACKEND\_PORT}      & \texttt{4000}       & Puerto del backend Elixir \\
\texttt{FRONTEND\_PORT}     & \texttt{3000}       & Puerto del frontend Next.js \\
\texttt{AI\_ENGINE\_PORT}   & \texttt{8000}       & Puerto del motor de IA \\
\hline
\end{tabular}
\end{table}

\section{Primera puesta en marcha}

La primera vez es necesario ejecutar dos comandos en orden. El primero prepara el entorno: detiene servicios previos, construye las imágenes Docker, crea y migra las bases de datos y carga los datos iniciales:

\begin{lstlisting}[language=bash]
make setup
\end{lstlisting}

Este proceso puede tardar varios minutos la primera vez, principalmente por la descarga de imágenes base y la compilación de dependencias Elixir. \textbf{No levanta los servicios}; una vez completado, hay que arrancarlos explícitamente:

\begin{lstlisting}[language=bash]
make up
\end{lstlisting}

\section{Arranque y parada habituales}

A partir de la segunda vez, el ciclo de trabajo diario se reduce a dos comandos:

\begin{lstlisting}[language=bash]
make up     # Levanta todos los servicios en segundo plano
make down   # Detiene todos los servicios
\end{lstlisting}

Tras ejecutar \texttt{make up} los servicios quedan accesibles en las siguientes URLs (tabla~\ref{tab:urls}).

\begin{table}[h]
\centering
\caption{URLs de acceso a los servicios en entorno local.}
\label{tab:urls}
\begin{tabular}{ll}
\hline
\textbf{Servicio} & \textbf{URL} \\
\hline
Frontend (Next.js)    & \url{http://localhost:3000} \\
Backend API (Phoenix) & \url{http://localhost:4000} \\
Motor de IA (FastAPI) & \url{http://localhost:8000} \\
Backoffice admin (Phoenix) & \url{http://localhost:4000/admin} \\
Mailpit (bandeja SMTP de pruebas) & \url{http://localhost:8025} \\
PostgreSQL            & \texttt{localhost:5432} \\
\hline
\end{tabular}
\end{table}

\section{Modos de ejecución}

El sistema admite tres modos de ejecución en función de la disponibilidad de hardware y del estado del modelo entrenado (tabla~\ref{tab:modos}).

\begin{table}[h]
\centering
\caption{Modos de ejecución disponibles.}
\label{tab:modos}
\begin{tabular}{lll}
\hline
\textbf{Modo} & \textbf{Comando} & \textbf{Requisito} \\
\hline
GPU  & \texttt{make up}      & GPU NVIDIA + modelo entrenado \\
CPU  & \texttt{make cpu-up}  & Modelo entrenado (inferencia lenta) \\
Mock & \texttt{make mock-up} & Ninguno (respuestas simuladas) \\
\hline
\end{tabular}
\end{table}

El modo \textit{mock} es el recomendado para desarrollo frontend y pruebas funcionales sin necesidad de disponer del modelo entrenado.

\section{Usuarios de prueba}

El comando \texttt{make setup} (a través de \texttt{mix run priv/repo/seeds.exs}) crea automáticamente los siguientes usuarios con la cuenta ya confirmada (tabla~\ref{tab:seeds}).

\begin{table}[h]
\centering
\caption{Usuarios de prueba creados por el seed.}
\label{tab:seeds}
\begin{tabular}{llll}
\hline
\textbf{Email} & \textbf{Contraseña} & \textbf{Username} & \textbf{Rol} \\
\hline
\texttt{admin@test.com} & \texttt{password123} & \texttt{admin} & Administrador \\
\hline
\end{tabular}
\end{table}

El usuario administrador tiene acceso al backoffice en \url{http://localhost:4000/admin}.

\section{Consulta de logs}

Para inspeccionar la salida de cualquier servicio en tiempo real:

\begin{lstlisting}[language=bash]
make logs   # Menú interactivo para seleccionar el servicio
\end{lstlisting}

\section{Reset de base de datos}

Si es necesario partir de un estado limpio sin reconstruir imágenes:

\begin{lstlisting}[language=bash]
make db-reset   # Drop + migrate + seed
\end{lstlisting}
